\scanpage{66}
Supposons en particulier que $X$ \unsure{soit} \uline{lisse} sur $S$, (donc on peut prendre $Y=X$, mais peu importe...) et $S_0\subset S$ une imm. nilpotente, d'où $X_0=X\times_S S_0\longrightarrow X$ imm. nilpotente. Considérons $X_0$ comme sch \uline{sur $S$} (pas sur $S_0$) pour définir $X_{0\,\mathrm{cris}}$.
On trouve alors
\[
H^*(X_{\mathrm{cris}},\mathcal O_{X_{\mathrm{cris}}})=H^*(X_{0\,\mathrm{cris}},\mathcal O_{X_{0\,\mathrm{cris}}})
\]
\uline{Cette relation prouve que le} 1\textsuperscript{er} \uline{mb ne dépend pas} (mod isom canonique) \uline{de la façon dont on a relevé le schéma lisse} $X_0/S_0$ \uline{en un schéma lisse sur} $S$.
Lorsque $S$ est de car. nulle, on verra que le premier mb est isom\inserted{(can. isom :)} à $H^*(X_{\mathrm{zar}},\Omega^{\boldsymbol\cdot}_{X/S})$, et on aura donc trouvé une définition des \uline{\unsure{stratifications}} \uline{can} \uline{en coh. de De Rham}\inserted{pour un morphisme lisse \deleted{\unsure{$f:X\longrightarrow S$}} \unsure{quasi-proj} $f:X\longrightarrow S$,} sur une base $S$ de car. nulle.
(Même mieux, on \deleted{\unsure{pourra en un}} \unsure{saura} arriver à définir $Rf_*(\Omega^{\boldsymbol\cdot}_{X/S})$ comme un \uline{cristal absolu} sur $S$ .....\inserted{ainsi l'étendre à \uline{tt} voisinage infinitésimal de $S$, pas seulement ceux qui loc. se rétractent sur $S$}).
\mspage{65}
